\section{Week 3 API Memory Reachability}
\label{sec:week3-api-memory}

\paragraph{Research question and claim.}
The strong secondary question is whether the packed-key relation used in the
HAETAE KeyGen/Sign algorithms (official PDF, Figure~8) survives the actual
Jasmin export/import loops.  The whole claim is tracked as
\claimid{OBL-API-KEY-MEMORY-RAW}:
\[
  \mathsf{KeyGen.local}\;(vk,sk)
  \Longrightarrow
  \mathsf{memory}[vku,vku+992)=vk
  \land
  \mathsf{memory}[sku,sku+1408)=sk
  \Longrightarrow
  \prefix_{992}(sk_{\mathrm{Sign}})=vk_{\mathrm{Verify}}.
\]

\paragraph{Pinned implementation and extraction.}
Theory \theoryname{ApiKeyMemoryBridge} opens the generated procedures
\procname{__kp_api_copy_2080_to_addr} and
\procname{__kp_api_copy_2752_to_addr} from \pathname{keypair.jazz}, plus the
Sign and Verify copies of \procname{_api_copy_raw_to_2752_prefix}.  The source
SHA-256 values are \texttt{8b2ddac2...3e01} for KeyGen,
\texttt{fd58271f...ac31} for Sign, \texttt{789f36e7...b06} for Verify, and
\texttt{4d04b8e7...a8a0} for their shared \pathname{api.jinc}.  Regenerated
theory hashes are recorded in
\pathname{manifests/generated-extractions.sha256}; regeneration is part of
the main verifier.

\paragraph{Compiled generated-helper theorems.}
The following actual procedure-level results fresh-compile:
\begin{itemize}
  \item \theoryname{keygen_export_vk_mode2_prefix}: under
        \theoryname{valid_region_w64} and length 992, the VK exporter writes
        every byte in the external interval to the corresponding local VK
        byte;
  \item \theoryname{keygen_export_sk_mode2_prefix}: the analogous statement
        for all 1408 mode-2 SK bytes;
  \item \theoryname{keygen_export_sk_mode2_prefix_frames_vk}: under the
        explicit minimal \theoryname{disjoint_regions} premise, the same
        actual SK exporter also preserves a previously established 992-byte
        VK memory prefix;
  \item \theoryname{sign_import_mode2_sk_prefix}: under
        \theoryname{valid_region_int} and length 1408, the actual Sign
        importer returns a local array whose first 992 bytes equal memory;
  \item \theoryname{verify_import_mode2_vk_prefix}: the actual Verify
        importer returns the 992-byte VK prefix; and
  \item \theoryname{sign_verify_api_importer_same_inputs}: identical memory,
        source, and length inputs produce identical results for the two
        generated importer copies.
\end{itemize}
These are Hoare partial-correctness or exact equivalence statements.  They do
not assume or conclude public API termination beyond the finite extracted
loops, nor do they manufacture a desired memory with an existential store.

\paragraph{Proof strategy and loop invariant.}
Each exporter invariant records the processed counter interval and a
\theoryname{mem2080_prefix} or \theoryname{mem2752_prefix} relation.  The
no-wrap premise justifies converting the generated word addition into integer
address addition.  Each importer first proves equality for indices below its
requested length, then shows the zero-fill tail cannot modify indices below
992.  This last frame fact is essential for Sign, whose mode-2 SK length is
1408 although the downstream hash reads only its 992-byte VK prefix.

\paragraph{Adapters into the mu theorem.}
The pure, compiled lemmas
\theoryname{exported_mode2_regions_agree},
\theoryname{imported_prefixes_agree}, and
\theoryname{imported_sign_sk_reaches_mu_memory} compose the packed-key prefix
with the helper postconditions and obtain exactly
\theoryname{ExtractedMuHashPrefix.sk_memory_prefix}.  Thus the dependency
graph has reached
\[
  \prefix_{992}(sk)=vk
  \rightarrow \mathsf{actual\ exports}
  \rightarrow \mathsf{actual\ imports}
  \rightarrow \mathsf{sk\_memory\_prefix},
\]
but the arrows are not yet a single sequential public-call theorem.

\paragraph{Address-model difference and aliasing.}
The generated KeyGen exporters retain \texttt{W64.t} addresses.  In contrast,
the selected Sign/Verify helper parameters originate as Jasmin \texttt{ui64}
but are emitted by \texttt{jasmin2ec} as mathematical \texttt{int}.  The
theory therefore keeps \theoryname{valid_region_w64} and
\theoryname{valid_region_int} distinct; no axiom equates their addresses.
Concrete zero/adjacent addresses witness both predicates and
\theoryname{disjoint_regions}, so the contracts are satisfiable.

Separation is correctness-critical when the later SK export must preserve the
earlier VK export.  It is not required merely to compare two isolated,
read-only importer calls.  A stronger all-buffer separation policy would be a
proof convenience and has not been adopted.  Time-dependent aliasing at the
unextracted public callers remains unanalyzed.

\paragraph{Failed boundary and next minimum obligation.}
The first attempt to concatenate the two exporters in a proof-only pair
wrapper became unstable at the intermediate global-memory ghost binding.
Replacing that wrapper with a stronger postcondition on the actual SK
exporter succeeded: the compiled frame theorem retains the earlier VK
predicate bytewise.  What remains cannot be solved by that local frame alone.
The word and integer pointers still cannot be identified syntactically, and
adding such an identification as an axiom was rejected.  The next minimum
theorem must compose the two actual exporter calls at their real caller, bind
\(\mathsf{to\_uint}(sku),\mathsf{to\_uint}(vku)\) to the two importer integer
addresses, and carry the resulting memory through the public call sequence.

\paragraph{Verification and status.}
The focused extraction script and
\texttt{easycrypt compile -script -no-eco} compile all results above; the
combined command is \pathname{scripts/verify-all.sh}.  The generated helper
claims are \status{PROVED} as partial correctness.  The final run reports 14
authored targets compiled with \texttt{-no-eco}, 20 selected baselines, a
successful notes build, and unchanged read-only roots.  The whole
\claimid{OBL-API-KEY-MEMORY-RAW} is \status{PARTIAL}, with the exact residual
being actual-caller export composition, address binding, and public-call
orchestration.
